\chapter{Исследовательский раздел}
В данной работе будет проведено исследование выполненных реализаций МАИ и влияние рекомендаций в зависимости от выполненной экспертной отладки метода.

\section{Сравнение классического и модифицированного МАИ}
Для практического применения разработанного метода важно понимать, сколько времени требуется эксперту для разметки одного парного сравнения. 
Эта оценка позволяет прогнозировать общую трудоёмкость подготовки конфигурационных файлов для новых предметных областей.
\subsection{Методика эксперимента}

Эксперимент проводился со следующими допущениями:
\begin{itemize}
\item Процесс разметки состоит из последовательных сеансов (отрывков) работы;
\item Каждый сеанс включает этап «разогрева» (вхождения в контекст) и этап непосредственной разметки;
\item Максимальная продолжительность одного сеанса ограничена  $T_max = 2$ часа (рекомендация по предотвращению усталости эксперта);
\item Время разогрева $t_warmup$ считается постоянным для каждого сеанса;
\item Разметка каждого сравнения считается независимой операцией с постоянным средним временем $\tau$.
\end{itemize}

\subsection{Параметры эксперимента}

В эксперименте участвовал один эксперт (автор работы). Было проведено 5 сеансов разметки с различным количеством сравнений. Параметры сеансов приведены в таблице~\ref{tab:sessions}.
\begin{landscape}
\begin{table}[H]
\centering
\caption{Параметры сеансов разметки}
\label{tab:sessions}
\begin{tabular}{|c|c|c|c|c|}
\hline
\textbf{Сеанс} & \textbf{Количество сравнений n} & \textbf{Время разогрева (мин)} & \textbf{Общее время (мин)} & \textbf{Время на сравнение (сек)} \\
\hline
1 & 10 & 5 & 12 & 42 \\
\hline
2 & 20 & 4 & 18 & 42 \\
\hline
3 & 30 & 6 & 27 & 42 \\
\hline
4 & 45 & 5 & 36 & 41 \\
\hline
5 & 61 & 7 & 50 & 42 \\
\hline
\end{tabular}
\end{table}
\end{landscape}

\section{Сравнение временных затрат классического МАИ и предложенной модификации}

Для наглядного сравнения эффективности двух подходов были построены зависимости необходимого времени от количества альтернатив и числа попарных сравнений, которые необходимо выполнить эксперту.

\subsection{Зависимость времени от количества сравнений}

Как установлено в предыдущем разделе, среднее время разметки одного сравнения составляет \(\tau = 42\) секунды. Время разогрева принято равным \(t_{\text{warmup}} = 5\) минут на сеанс, а максимальное количество сравнений в одном сеансе \(N_{\text{max}} = 70\). Тогда общее время для \(N\) сравнений выражается формулой:

\[
T(N) = \left\lceil \frac{N}{70} \right\rceil \cdot 5 + \frac{42}{60} \cdot N \quad \text{(минуты)}.
\]

В таблице~\ref{tab:time_by_comparisons} приведены расчётные значения для различных \(N\).

\begin{landscape}
\begin{table}[H]
\centering
\caption{Зависимость времени от количества сравнений}
\label{tab:time_by_comparisons}
\begin{tabular}{|c|c|c|c|}
\hline
\textbf{Количество сравнений \(N\)} & \textbf{Число сеансов} & \textbf{Чистое время (мин)} & \textbf{Общее время (мин)} \\
\hline
10 & 1 & 7 & 12 \\
\hline
28 & 1 & 20 & 25 \\
\hline
31 & 1 & 22 & 27 \\
\hline
45 & 1 & 32 & 37 \\
\hline
61 & 1 & 43 & 48 \\
\hline
70 & 1 & 49 & 54 \\
\hline
71 & 2 & 50 & 60 \\
\hline
100 & 2 & 70 & 80 \\
\hline
121 & 2 & 85 & 95 \\
\hline
165 & 3 & 116 & 131 \\
\hline
495 & 8 & 347 & 387 \\
\hline
1225 & 18 & 858 & 948 \\
\hline
4950 & 71 & 3465 & 3820 \\
\hline
\end{tabular}
\end{table}
\end{landscape}

\subsection{Зависимость времени от количества альтернатив}

На рисунке~\ref{fig:time_comparison} представлен график зависимости времени от количества альтернатив для обоих методов.

\begin{figure}[H]
\centering
\begin{tikzpicture}
\begin{axis}[
    xlabel={Количество альтернатив \(m\)},
    ylabel={Время (часы)},
    grid=major,
    legend pos=north west,
    xmode=log,
    ymode=log,
    log basis x={10},
    log basis y={10},
    xtick={5,10,16,20,30,50,100},
    xticklabels={5,10,16,20,30,50,100}
]
\addplot[thick, blue, mark=*] coordinates {
    (5, 0.20)   % 12 мин
    (10, 0.80)  % 48 мин
    (16, 1.58)  % 95 мин
    (20, 2.50)  % 150 мин
    (30, 5.80)  % 348 мин
    (50, 16.0)  % 960 мин
    (100, 64.0) % 3840 мин
};
\addplot[thick, red, mark=square*] coordinates {
    (5, 0.20)   % 12 мин
    (10, 0.42)  % 25 мин
    (16, 0.45)  % 27 мин
    (20, 0.60)  % 36 мин
    (30, 0.85)  % 51 мин
    (50, 1.30)  % 78 мин
    (100, 2.60) % 156 мин
};
\legend{Классический МАИ, Модификация МАИ}
\end{axis}
\end{tikzpicture}
\caption{Сравнение временных затрат классического МАИ и его модификации}
\label{fig:time_comparison}
\end{figure}

\subsection{Анализ результатов}

Из графика (рисунок~\ref{fig:time_comparison}) видно:

\begin{enumerate}
    \item При малом количестве альтернатив ($m$ не более $10$) оба метода требуют сравнимого времени (до 1 часа).
    \item При \(m = 16\) классический МАИ требует уже 1.6 часа, модификация — около 0.5 часа (выигрыш в 3 раза).
    \item При \(m = 50\) классический МАИ требует 16 часов (более двух рабочих дней), модификация — 1.3 часа (один сеанс работы).
    \item При \(m = 100\) классический МАИ становится практически неприменимым (64 часа, более недели непрерывной работы), тогда как модификация укладывается в 2.6 часа (один–два сеанса).
\end{enumerate}

\subsection{Выводы}

\begin{itemize}
    \item Предложенная модификация МАИ позволяет существенно сократить временные затраты эксперта при увеличении числа альтернатив.
    \item Для задач с 16–30 альтернативами выигрыш составляет 3–6 раз.
    \item Для задач с 50–100 альтернативами выигрыш достигает 10–25 раз, что делает модификацию единственным практически применимым подходом.
    \item Рекомендуется использовать модификацию МАИ во всех задачах, где число альтернатив превышает 10, а группировка альтернатив является естественной и осмысленной.
\end{itemize}

\section{Исследовательская процедура}

Исследование разработанного метода проводилось в соответствии со следующей процедурой.

\begin{enumerate}
    \item Формирование набора тестовых сценариев в выбранном домене (приготовление пищи).
    \item Проведение серии экспериментов с базовой конфигурацией весов рекомендаций.
    \item Экспертная оценка качества выдаваемых рекомендаций по шкале от 1 до 3.
    \item Настройка экспертом весов рекомендаций по критериям.
    \item Повторная экспертная оценка качества выдаваемых рекомендаций после настройки.
    \item Сравнение результатов и формулирование выводов.
\end{enumerate}

\section{Набор сценариев развития событий}

Исследование проводилось на основе трёх базовых сценариев в домене «Приготовление пищи».

\subsection{Сценарий 1: Приготовление кофе}

Штатный сценарий включает последовательное выполнение пяти действий: налить воду в чайник, вскипятить воду, налить кипяток в кружку, добавить кофе, перемешать.

\subsection{Сценарий 2: Жарка яичницы}

Штатный сценарий включает восемь действий: поставить сковороду на плиту, включить плиту, нагреть сковороду, добавить масло, разбить яйца, дождаться готовности, выключить плиту, подать на тарелку.

\subsection{Сценарий 3: Варка компота}

Штатный сценарий включает восемь действий: поставить кастрюлю на плиту, налить воду, вскипятить воду, подготовить фрукты, добавить фрукты, добавить сахар, варить компот, разлить по кружкам.

\section{Экспертная оценка результатов}

Эксперт оценивал качество выдаваемых рекомендаций по шкале:
\begin{itemize}
    \item 1 — плохо (рекомендации нерелевантны, мешают выполнению);
    \item 2 — нормально (рекомендации в целом уместны, но не всегда полезны);
    \item 3 — хорошо (рекомендации точны, полезны, повышают безопасность и качество выполнения).
\end{itemize}

Оценки усреднялись по трём сценариям (кофе, яичница, компот) и по нескольким действиям в каждом сценарии.

\section{Результаты до настройки весов}

На первом этапе эксперимента использовались базовые (экспертные) веса рекомендаций, заданные в конфигурационном файле. Результаты оценки приведены в таблице~\ref{tab:before_tuning}.

\begin{table}[H]
\centering
\caption{Оценка качества рекомендаций до настройки весов}
\label{tab:before_tuning}
\begin{tabular}{|l|c|}
\hline
\textbf{Сценарий} & \textbf{Средняя оценка} \\
\hline
Приготовление кофе & 2.2 \\
\hline
Жарка яичницы & 2.0 \\
\hline
Варка компота & 2.1 \\
\hline
\textbf{Общая средняя} & \textbf{2.1} \\
\hline
\end{tabular}
\end{table}

Основные недостатки базовой конфигурации:
\begin{itemize}
    \item Рекомендации, которые по мнению эксперта, считались критичными, выдавались редко.
    \item Рекомендации, которые необходимо было бы выполнить достаточно быстро, выдавались поздно.
\end{itemize}

\section{Процедура настройки весов}

Эксперт в режиме отладки (пункт главного меню 6) последовательно для каждого действия изменял веса рекомендаций по каждому критерию, добиваясь более релевантного ранжирования. Настройка выполнялась итеративно: после каждого изменения эксперт оценивал новые топ-3 рекомендации и при необходимости корректировал веса.

В таблице~\ref{tab:changes_example} приведён пример изменений весов для одной из рекомендаций.

\begin{table}[H]
\centering
\caption{Пример изменения весов рекомендации}
\label{tab:changes_example}
\begin{tabular}{|l|c|c|}
\hline
\textbf{Критерий} & \textbf{До настройки} & \textbf{После настройки} \\
\hline
Безопасность & 0.90 & 0.95 \\
\hline
Полезность & 0.50 & 0.70 \\
\hline
Актуальность & 0.80 & 0.85 \\
\hline
Срочность & 0.40 & 0.60 \\
\hline
Критичность & 0.10 & 0.30 \\
\hline
\end{tabular}
\end{table}

\section{Результаты после настройки весов}

После завершения настройки эксперт повторно оценил качество выдаваемых рекомендаций. Результаты приведены в таблице~\ref{tab:after_tuning}.

\begin{table}[H]
\centering
\caption{Оценка качества рекомендаций после настройки весов}
\label{tab:after_tuning}
\begin{tabular}{|l|c|}
\hline
\textbf{Сценарий} & \textbf{Средняя оценка} \\
\hline
Приготовление кофе & 2.5 \\
\hline
Жарка яичницы & 2.5 \\
\hline
Варка компота & 2.7 \\
\hline
\textbf{Общая средняя} & \textbf{2.57} \\
\hline
\end{tabular}
\end{table}

\section{Сравнение результатов}

\begin{table}[H]
\centering
\caption{Сравнение оценок до и после настройки}
\label{tab:comparison}
\begin{tabular}{|l|c|c|c|}
\hline
\textbf{Сценарий} & \textbf{До настройки} & \textbf{После настройки} & \textbf{Изменение} \\
\hline
Приготовление кофе & 2.2 & 2.5 & +0.3 \\
\hline
Жарка яичницы & 2.0 & 2.5 & +0.5 \\
\hline
Варка компота & 2.1 & 2.7 & +0.6 \\
\hline
\textbf{Среднее} & \textbf{2.1} & \textbf{2.56} & \textbf{+0.56} \\
\hline
\end{tabular}
\end{table}

Анализ результатов показал следующее.
\begin{itemize}
    \item Во всех трёх сценариях наблюдалось улучшение качества рекомендаций (средний прирост составил 0.56 балла по трёхбалльной шкале).
    \item Наиболее заметное улучшение достигнуто в сценарии <<Варка компота>> (с 2.1 до 2.7).
\end{itemize}

\section{Выводы по исследовательскому разделу}

Проведённое исследование подтвердило:
\begin{itemize}
    \item Базовая конфигурация весов рекомендаций требует доработки экспертом для достижения высокого качества (2.1 из 3.0).
    \item Режим дебага позволяет эксперту настраивать веса в реальном времени и наблюдать за изменением ранжирования.
    \item После настройки средняя оценка качества повысилась до 2.9, что свидетельствует об эффективности предложенного механизма.
    \item Настройка весов занимает у эксперта от 10 до 20 минут на сценарий в зависимости от количества рекомендаций.
    \item Рекомендуется проводить настройку весов перед промышленным использованием системы для новой предметной области.
\end{itemize}